\section{Terminology}
\label{sec:terminology}

The vocabulary extends established terms with \emph{security context}.

\paragraph{Properties and assets.}
A \emph{security property} is an attribute of an asset that the system must
preserve: confidentiality, integrity and availability~\citep{avizienis2004},
refined by ISO/IEC~25010 into confidentiality, integrity, non-repudiation,
accountability, authenticity and resistance~\citep{iso25010}. An \emph{asset} is
data or functionality whose security properties matter; for a persisted
Highscore board the assets are the record store and the displayed ranking.

\paragraph{Threats and attack surface.}
A \emph{threat} is a potential violation of a security property; an
\emph{attack} is an intentional attempt to realise it~\citep{rfc4949}.
CAPEC catalogues \emph{attack patterns}~\citep{capec}; STRIDE guides threat
enumeration per system element~\citep{shostack2014}. A
\emph{trust boundary} separates components with different levels of
trust~\citep{shostack2014}; the set of entry points, exit points and channels
through which untrusted data reaches the system is its \emph{attack
surface}~\citep{manadhata2011}. A \emph{threat model} fixes which actors and
channels are in scope; here it is a desktop game whose untrusted inputs are
local files and content from a shared level server.

\paragraph{Weaknesses and vulnerabilities.}
A \emph{weakness} is a type of flaw that can lead to a violation of a security
property; a \emph{vulnerability} is an exploitable occurrence of a weakness in
a specific system~\citep{cwe,cve}. CWE names weakness types~\citep{cwe}; the
Seven Pernicious Kingdoms give a coarser taxonomy~\citep{tsipenyuk2005}. In
the fault--error--failure chain, a weakness is a fault and an exploited
vulnerability is a security failure~\citep{avizienis2004}. The executable
issue checks detect weaknesses; they do not establish vulnerabilities.

\paragraph{Requirements, controls, verification.}
A \emph{security requirement} is a constraint on a functional requirement that
protects an asset against a threat~\citep{haley2008,firesmith2003}; misuse
cases are one elicitation technique~\citep{sindre2005}. A \emph{security
control} is a safeguard that satisfies a security requirement~\citep{nist80053}.
At code level, security controls follow design principles~\citep{saltzer1975}
and secure coding rules~\citep{long2011}. Its \emph{enforcement point} is the
operation at which the security control must hold. \emph{Verification} checks
that a security control satisfies its security requirement; ASVS provides
verification requirements per security
area~\citep{asvs5}, and secure development frameworks place them in the
lifecycle~\citep{nist800218}.

\paragraph{Security context.}
For a repository $R$ and a change $\Delta$, the \emph{security context}
$\mathrm{SC}(R,\Delta)$ is the set of typed statements supplied to the
generator before generation, each grounded in inspected source or explicitly
marked as prospective: the assets and entry points $\Delta$ touches, the trust
boundaries it crosses, the applicable weakness types (CWE), the derived
security requirements with enforcement points, proposed security controls,
verification hints, and unknowns. Security context is distinct from
\emph{security feedback}, the outcome of tests on generated code, which is
never returned to the generator, and from generic secure coding instructions
in the prompt~\citep{tony2025}, which contain no repository evidence.

\paragraph{Mapping to this study.}
Table~\ref{tab:mapping} places the evaluation categories in this vocabulary.
Each issue check probes one weakness type through a fixed fixture; the
categories together cover the input, retention, parsing, resource and
deserialization behaviour of one feature and are not a vulnerability
inventory.

\begin{table}[t]
\centering
\footnotesize
\setlength{\tabcolsep}{4pt}
\caption{Evaluation categories of the Highscore study in standard terms.}
\label{tab:mapping}
\begin{tabular}{@{}llll@{}}
\toprule
Category & Weakness type & Property & Requirement source \\
\midrule
Input policy & CWE-20, CWE-1284 & Integrity & Contract; ASVS validation; CERT IDS \\
Retention policy & CWE-770 & Availability & Contract \\
Parser robustness & CWE-755, CWE-20 & Availability & CERT ERR \\
Resource stress & CWE-400, CWE-770, CWE-789 & Availability & CERT MSC, FIO \\
Deserialization dispatch & CWE-502 & Integrity & CERT SER; ASVS \\
\bottomrule
\end{tabular}
\end{table}
